\section{Introduction}\label{sec:intro}
ThermalSense is an interactive OpenGL~3.3 (core profile) base-defence game with a full day--night cycle, written in C++17 with GLFW and GLAD and no
model or image files. Every visible solid is a \emph{part}: one of the lab's unit primitives (cube, sphere, cylinder) drawn with its own
$M=T\,R\,S$ matrix. Two UAVs patrol the base, a computer-vision loop tracks heat sources, and the player or the squad defends the site against intruding rovers. Version~6 adds the sun and the sky (Chapter~\ref{sec:daycycle}), a garrison of \nSoldiers{} animated soldiers (Chapter~\ref{sec:soldiers}), seeded random intruder routes (Chapter~\ref{sec:randompaths}) and a modelled thermal sensor whose noisy images are the only way to a kill (Chapter~\ref{sec:realcv}).

\paragraph{Requirements.} The lab sheet (\texttt{graphics-lab-04.pdf}) asks for a canonical view volume; unit primitives centred at the origin
(36 slices $=10^\circ$); counter-clockwise front faces; indexed drawing (12 floats $+$ 6 indices instead of 18 floats) with an explicit VAO/VBO/EBO layout;
a sphere built from one reflected hemisphere; a cylinder along $z$; an $S,R,T$ table per component with $M=T\cdot R\cdot S$; and every object built from primitives only,
documented in an Excel table (T01--T14). The project proposal adds three shading models (Flat, Gouraud, Phong), three light sources, a hierarchical searchlight, a Z-buffer,
wireframe modes, camera controls and a thermal colour map. Table~\ref{tab:course} lists the course-level deliverables and where they are met.

\begin{table}[h]
\centering\small
\caption{Course checklist and where it is met. ``Pending'' items are separate deliverables.}
\label{tab:course}
\begin{tabular}{@{}p{6.1cm}p{9cm}l@{}}
\toprule
Item & Where / how & Status\\\midrule
Dynamic and controllable objects & \nPossessable{} possessable objects (machines, soldiers, SKY clock), Section~\ref{sec:control} & \sDone\\
Complex motion (v6) & Spline routes with arc length, gait by FK + two-link IK, sun path (Sections~\ref{sec:daycycle}, \ref{sec:soldiers}) & \sDone\\
Complex motion & Track links, spring-damper attitude, rotor cyclic, Gatling spin-up (Section~\ref{sec:control}) & \sDone\\
Report with an audit & This document, Appendix~\ref{app:audit} & \sDone\\
Slide presentation (10 slides, outline on slide 2) & \code{ThermalSense\_Presentation.pptx} & \sDone\\
2-minute video with sound & \code{demo\_v5.mp4}, embedded on the last slide; the v6 cut is not recorded yet & \sPart\\
Excel workbook & generated from \texttt{scene\_export.json}, verified; crops in this report & \sDone\\\bottomrule
\end{tabular}
\end{table}

\paragraph{The site.}
The scene is a base of 36\,m radius (terrain mesh) with a 34\,m HESCO perimeter, surrounded by hills and a tree line. In the middle lies the 20\,m patrol compound where the UAVs and the three intruder
rovers play hide-and-seek, with the command bunker buried under it. Figure~\ref{fig:plan} shows the site plan from \code{src/objects/BaseLayout.h}. All coordinates are world metres, $+Z$ north, $+X$ east, $y$ up.

\begin{figure}[h]
\centering
\begin{tikzpicture}[scale=0.118,font=\scriptsize]
  \draw[fill=black!4] (0,0) circle (36);
  \draw[dashed] (0,0) circle (34);
  \draw[fill=green!8,thin] (0,0) circle (20);
  \node[above] at (0,20.5) {compound r\,=\,20};
  % bunker footprint
  \draw[fill=gray!35] (-5.5,-8.5) rectangle (5.5,4);
  \node at (0,-2) {bunker (underground)};
  % ramp
  \draw[fill=gray!60] (-1.5,6) rectangle (1.5,13);
  \node[right] at (1.6,9.5) {ramp};
  % road
  \draw[thick] (-2.2,19) -- (-2.2,34) (2.2,19) -- (2.2,34);
  \node[above] at (0,34.3) {outer gate};
  % helipad
  \draw[fill=blue!15] (20.3,17) circle (5); \node at (20.3,17) {helipad};
  % radar
  \draw[fill=orange!30] (29,-3) circle (2.7); \node[below=2mm] at (29,-3) {radar};
  % AA gun
  \draw[fill=red!20] (21,-19) circle (3.5); \node[below=3mm] at (21,-19) {AA gun};
  % motor pool
  \draw[fill=olive!30,rotate around={45:(-22,-16)}] (-25,-19) rectangle (-19,-13); \node[below=4mm] at (-22,-16) {motor pool/tank};
  % barracks
  \draw[fill=brown!25] (-27,11) circle (4); \node at (-27,11) {barracks};
  % depot
  \draw[fill=yellow!25] (-6,-28.5) circle (3); \node[right=2mm] at (-6,-28.5) {fuel depot};
  % flag
  \fill (-4.5,27) circle (0.8); \node[left=1mm] at (-4.5,27) {flag};
  % towers
  \foreach \a in {60,130,245,340} {\fill[red!60!black] ({32.4*cos(\a)},{32.4*sin(\a)}) circle (1.1);}
  % north arrow
  \draw[-{Latex}] (-34,28) -- (-34,34) node[above] {$+Z$};
  \draw[-{Latex}] (-34,28) -- (-28,28) node[right] {$+X$};
\end{tikzpicture}
\caption{Plan view of the site (from \code{BaseLayout.h}; tower angles are measured from $+X$ towards $+Z$ as in
\code{Base.cpp}). Dark circles: the four guard towers at $r=32.4$\,m.}
\label{fig:plan}
\end{figure}

\subsection{The journey of one frame}
The report follows the data. A \emph{vertex} begins as a number in a unit primitive; the $T\,R\,S$ matrix of its part and the hierarchy above it place it in the world (Chapters~\ref{sec:builder}--\ref{sec:construction}); the view and projection matrices, clipping and the viewport turn
it into a pixel and a depth (Chapter~\ref{sec:geomchain}); the rasteriser interpolates its attributes (Chapter~\ref{sec:interp}); the fragment shader lights it (Chapter~\ref{sec:shading}); the reflection and post passes finish the image (Chapters~\ref{sec:reflection},~\ref{sec:pipeline}).
The scene is then rendered a second time as heat for each UAV; the computer vision turns that image into tracks (Chapter~\ref{sec:vision}) and the squad and rovers act on them (Chapters~\ref{sec:rover}, \ref{sec:squad}), which changes the next frame. A strip at the head of each
chapter marks its place in this pipeline; workbook screenshots accompany the objects they document.

\subsection{System overview}
Figure~\ref{fig:sysov}: the simulation updates a scene graph of primitives; the renderer draws it through the shading and post chain, and renders it again from each UAV's gimbal as a heat image for the vision. Audio and HUD read the game events.
\begin{figure}[h]
\centering
\begin{tikzpicture}[node distance=5mm and 7mm, every node/.style={font=\scriptsize},
  box/.style={draw,rounded corners=2pt,minimum height=8mm,minimum width=22mm,align=center,fill=blue!6},arr/.style={-{Latex[length=1.6mm]},thick}]
  \node[box] (prim) {unit primitives\\cube, sphere, cylinder};
  \node[box,right=of prim] (build) {Builder tables\\$M=TRS$, hierarchy};
  \node[box,right=of build] (scene) {scene graph\\\nLeaves{} parts};
  \node[box,below=of scene] (sim) {simulation\\controllables, rovers,\\combat, mission};
  \node[box,right=of scene,fill=orange!10] (rend) {renderer\\shading, shadows,\\SSR, post chain};
  \node[box,below=of rend,fill=orange!10] (heat) {heat render\\2 UAV cameras};
  \node[box,left=of sim,fill=green!8] (cv) {computer vision\\tracker, squad};
  \node[box,right=of rend] (screen) {screen\\+ HUD};
  \node[box,below=of screen,fill=purple!8] (audio) {audio engine\\procedural};
  \draw[arr] (prim)--(build); \draw[arr] (build)--(scene); \draw[arr] (scene)--(rend); \draw[arr] (rend)--(screen);
  \draw[arr] (scene)--(heat); \draw[arr] (heat)--(cv); \draw[arr] (cv)--(sim); \draw[arr] (sim)--(scene);
  \draw[arr] (sim.south) -- ++(0,-5mm) -| (audio.south); \draw[arr] (sim.south east) |- ([yshift=-3mm]screen.south west) ;
\end{tikzpicture}
\caption{System overview.}\label{fig:sysov}
\end{figure}

\paragraph{How to read this report.} Each chapter states the concept, the mathematics, the project's code (listings are lifted from the source by line range) and the tests or measurements.
Chapters~\ref{sec:prims}--\ref{sec:construction} cover geometry and object construction; \ref{sec:interp}--\ref{sec:pipeline} shading and the render pipeline; \ref{sec:rover} and \ref{sec:vision} the rover and the computer vision;
\ref{sec:v5plan}--\ref{sec:v5media} the game; the later chapters the machines, bunker, platform, tests and results; \ref{sec:daycycle}--\ref{sec:v6results} the v6 day cycle, soldiers, random routes, sensor model and their results. All numbers are generated from the program output.
